% Source: https://artofproblemsolving.com/community/c3h1557474_cube_problem
% Stable identity: aops-topic-1557474-q01
Let the large cube be $n\times n\times n$ (the figure shows the answer, $n=6$).

\begin{asy}
usepackage("amsmath");
size(4.5cm);
defaultpen(fontsize(10pt));
int n=6; pair k=(0.45,0.3);
for(int i=0;i<=n;++i){ draw((i,0)--(i,n),gray(0.55)); draw((0,i)--(n,i),gray(0.55));
  draw((n,i)--(n,i)+n*k,gray(0.55)); draw((n,0)+i*k--(n,n)+i*k,gray(0.55));
  draw((0,n)+i*k--(n,n)+i*k,gray(0.55)); draw((i,n)--(i,n)+n*k,gray(0.55)); }
draw((0,0)--(n,0)--(n,n)--(0,n)--cycle,linewidth(0.9));
draw((n,0)--(n,0)+n*k--(n,n)+n*k--(0,n)+n*k--(0,n),linewidth(0.9));draw((n,n)--(n,n)+n*k,linewidth(0.9));
label("$6$",(n/2,0),S);label("$6$",(0,n/2),W);label("$6$",(n,0)+n*k/2,SE);
label("\mbox{}",(min(currentpicture,user=true)+max(currentpicture,user=true))/2,black); // keeps the text after the figure black when LaTeX typesets the labels inline
\end{asy}

A unit cube is unpainted exactly when it does not touch the surface, i.e.\ it lies in the inner cube obtained by removing one layer from each face; that inner cube is $(n-2)\times(n-2)\times(n-2)$. So
\[
(n-2)^3=64\quad\Longrightarrow\quad n-2=4\quad\Longrightarrow\quad n=6,
\]
and the large cube is made of $6^3=\boxed{216}$ unit cubes. (The poster's answer $216$ is correct.)
